
\newpage
\section{Анализ необходимых исходных данных}
\label{part:data_intro}

Из описания предлагаемой методики анализа влияния моря на прибрежные территории становится очевидной критическая важность для расчетов доступных, точных и подробных данных.
Необходимый и минимальный перечень которых составляют:

\begin{itemize}
\item исторические метеорологические данные содержащие в себе информацию о силе и направлении ветра;
\item данные о батиметрии прибрежной зоны;
\item актуальное положение береговой линии.
\end{itemize}

При этом, не пренебрегая важностью прочих данных, без сомнения, наиболее интересными, как в плане получения, так и обработки являются метеоданные, принципиальным отличаем которых является их перспективная случайность и ретроспективная вариативность.
Принимая так же во внимание тот факт, что процедуры получения подобных данных могут представлять самодостаточную ценность для исследователей в смежных областях рассмотрению этого вопроса далее представляется первоочередное внимание.

